% Slug: ch20-merge
% Source: chapters/ch20.tex:150 (tikzpicture env start line)
% Original label: \label{fig:ch20-merge}
% Original caption (LaTeX source):
%   The join block sees both bindings of the pointer, so its target set is the union of the two branches.
% Original caption (plain text):
%   The join block sees both bindings of the pointer, so its target set is the union of the two branches.
% Required TikZ setup: \usepackage{tikz} + \usetikzlibrary{arrows.meta, positioning, shapes.geometric, fit} (loaded by support/compiler-pseudocode.sty).
% Shared styles used in this picture: `basic block', `cfg edge'.
%   Defined in support/compiler-pseudocode.sty: `basic block', `cfg edge', `back edge', `block label'.
% Usage: % Slug: ch20-merge
% Source: chapters/ch20.tex:150 (tikzpicture env start line)
% Original label: \label{fig:ch20-merge}
% Original caption (LaTeX source):
%   The join block sees both bindings of the pointer, so its target set is the union of the two branches.
% Original caption (plain text):
%   The join block sees both bindings of the pointer, so its target set is the union of the two branches.
% Required TikZ setup: \usepackage{tikz} + \usetikzlibrary{arrows.meta, positioning, shapes.geometric, fit} (loaded by support/compiler-pseudocode.sty).
% Shared styles used in this picture: `basic block', `cfg edge'.
%   Defined in support/compiler-pseudocode.sty: `basic block', `cfg edge', `back edge', `block label'.
% Usage: % Slug: ch20-merge
% Source: chapters/ch20.tex:150 (tikzpicture env start line)
% Original label: \label{fig:ch20-merge}
% Original caption (LaTeX source):
%   The join block sees both bindings of the pointer, so its target set is the union of the two branches.
% Original caption (plain text):
%   The join block sees both bindings of the pointer, so its target set is the union of the two branches.
% Required TikZ setup: \usepackage{tikz} + \usetikzlibrary{arrows.meta, positioning, shapes.geometric, fit} (loaded by support/compiler-pseudocode.sty).
% Shared styles used in this picture: `basic block', `cfg edge'.
%   Defined in support/compiler-pseudocode.sty: `basic block', `cfg edge', `back edge', `block label'.
% Usage: % Slug: ch20-merge
% Source: chapters/ch20.tex:150 (tikzpicture env start line)
% Original label: \label{fig:ch20-merge}
% Original caption (LaTeX source):
%   The join block sees both bindings of the pointer, so its target set is the union of the two branches.
% Original caption (plain text):
%   The join block sees both bindings of the pointer, so its target set is the union of the two branches.
% Required TikZ setup: \usepackage{tikz} + \usetikzlibrary{arrows.meta, positioning, shapes.geometric, fit} (loaded by support/compiler-pseudocode.sty).
% Shared styles used in this picture: `basic block', `cfg edge'.
%   Defined in support/compiler-pseudocode.sty: `basic block', `cfg edge', `back edge', `block label'.
% Usage: \input{graphs-code/ch20-merge.tikz.tex} inside a figure environment
%   (wrap with \begin{figure} ... \caption{...} \label{fig:ch20-merge} ... \end{figure}).
% Note: tikzpicture body below is verbatim from the source; do not reformat.
\begin{tikzpicture}[node distance=14mm and 20mm]
    \node[basic block] (ch20e) {E: p = \&x; q = \&y};
    \node[basic block,below left=of ch20e] (ch20t) {T: p = q};
    \node[basic block,below right=of ch20e] (ch20f) {F: };
    \node[basic block,below right=of ch20t] (ch20j) {J: r = p};
    \draw[cfg edge] (ch20e) -- (ch20t);
    \draw[cfg edge] (ch20e) -- (ch20f);
    \draw[cfg edge] (ch20t) -- (ch20j);
    \draw[cfg edge] (ch20f) -- (ch20j);
  \end{tikzpicture}
 inside a figure environment
%   (wrap with \begin{figure} ... \caption{...} \label{fig:ch20-merge} ... \end{figure}).
% Note: tikzpicture body below is verbatim from the source; do not reformat.
\begin{tikzpicture}[node distance=14mm and 20mm]
    \node[basic block] (ch20e) {E: p = \&x; q = \&y};
    \node[basic block,below left=of ch20e] (ch20t) {T: p = q};
    \node[basic block,below right=of ch20e] (ch20f) {F: };
    \node[basic block,below right=of ch20t] (ch20j) {J: r = p};
    \draw[cfg edge] (ch20e) -- (ch20t);
    \draw[cfg edge] (ch20e) -- (ch20f);
    \draw[cfg edge] (ch20t) -- (ch20j);
    \draw[cfg edge] (ch20f) -- (ch20j);
  \end{tikzpicture}
 inside a figure environment
%   (wrap with \begin{figure} ... \caption{...} \label{fig:ch20-merge} ... \end{figure}).
% Note: tikzpicture body below is verbatim from the source; do not reformat.
\begin{tikzpicture}[node distance=14mm and 20mm]
    \node[basic block] (ch20e) {E: p = \&x; q = \&y};
    \node[basic block,below left=of ch20e] (ch20t) {T: p = q};
    \node[basic block,below right=of ch20e] (ch20f) {F: };
    \node[basic block,below right=of ch20t] (ch20j) {J: r = p};
    \draw[cfg edge] (ch20e) -- (ch20t);
    \draw[cfg edge] (ch20e) -- (ch20f);
    \draw[cfg edge] (ch20t) -- (ch20j);
    \draw[cfg edge] (ch20f) -- (ch20j);
  \end{tikzpicture}
 inside a figure environment
%   (wrap with \begin{figure} ... \caption{...} \label{fig:ch20-merge} ... \end{figure}).
% Note: tikzpicture body below is verbatim from the source; do not reformat.
\begin{tikzpicture}[node distance=14mm and 20mm]
    \node[basic block] (ch20e) {E: p = \&x; q = \&y};
    \node[basic block,below left=of ch20e] (ch20t) {T: p = q};
    \node[basic block,below right=of ch20e] (ch20f) {F: };
    \node[basic block,below right=of ch20t] (ch20j) {J: r = p};
    \draw[cfg edge] (ch20e) -- (ch20t);
    \draw[cfg edge] (ch20e) -- (ch20f);
    \draw[cfg edge] (ch20t) -- (ch20j);
    \draw[cfg edge] (ch20f) -- (ch20j);
  \end{tikzpicture}
